\documentclass[a4paper,11pt]{article}
\usepackage[utf8]{inputenc}
\usepackage[T1]{fontenc}
\usepackage[french]{babel}
\usepackage[left=1cm,right=1cm,top=.5cm,bottom=0cm]{geometry}
\usepackage{enumitem}
\usepackage{hyperref}
\usepackage{xcolor}
\usepackage{titlesec}
\usepackage{lmodern}
\usepackage[normalem]{ulem}
\usepackage{pifont}

% --- Couleurs Crédit Agricole (vert) ---
\definecolor{primary}{RGB}{0, 100, 62}
\definecolor{secondary}{RGB}{80, 80, 80}

% --- Config Liens ---
\hypersetup{colorlinks=true, linkcolor=primary, filecolor=primary, urlcolor=primary}

% --- Titres ---
\titleformat{\section}{\large\bfseries\color{primary}\uppercase}{}{0em}{}[\titlerule]
\titlespacing{\section}{0pt}{8pt}{5pt}

% --- Commandes ---
\newcommand{\resumeItem}[1]{\item\small{#1}}

\begin{document}

% --- EN-TÊTE ---
\begin{center}
    {\Large \textbf{Loïc Aron MBASSI EWOLO}} \\ \vspace{4pt}
    \textbf{\large Développeur Back-End Java / API REST} \\
    \textit{\small Alternance 12 mois, disponible dès \textbf{septembre 2026}} \\ \vspace{4pt}
    \small
    \ding{168} Cergy, France (Mobile IDF) \quad
    \ding{37} (+33) 7 59 52 33 98 \quad
    \ding{41} \href{mailto:loicaron.mbassiewolo@ensea.fr}{loicaron.mbassiewolo@ensea.fr} \\ \vspace{2pt}
    \href{https://www.linkedin.com/in/loic-mbassi/}{LinkedIn : loic-mbassi} \quad
    \href{https://github.com/Nameless0l}{GitHub : Nameless0l} \quad
    \href{https://mbassiloic.tech}{Portfolio : mbassiloic.tech}
\end{center}

\vspace{-6pt}

% --- PROFIL ---
\noindent\small{Développeur backend avec 3 ans d'expérience en conception d'API REST, architectures microservices et systèmes de messaging asynchrone. Double diplôme \textbf{ENSEA} / \textbf{Polytechnique Yaoundé}, stage de recherche \textbf{INRIA Saclay} en cours. Expérience directe en \textbf{développement bancaire sécurisé} (Fintech) et en architecture réactive (\textbf{Spring Boot}, Redis, PostgreSQL). Je souhaite mettre ces compétences au service de la squad CRM Opérationnel de \textbf{CA-TS} pour contribuer à l'amélioration des outils de campagnes marketing et de gestion des leads.}

% --- FORMATION ---
\section{Formation}
\begin{itemize}[leftmargin=*, label=\tiny$\bullet$, nosep]
    \item \textbf{ENSEA -- École Nationale Supérieure de l'Électronique et de ses Applications} \hfill \textit{2025 -- 2027} \\
    \small{Double diplôme d'Ingénieur -- Systèmes Embarqués, Signal, IA. Cours : Architecture logicielle, Deep Learning, Traitement du Signal.}
    \item \textbf{École Polytechnique de Yaoundé (1\textsuperscript{ère} école d'ingénieur CEMAC)} \hfill \textit{2021 -- 2025} \\
    \small{Diplôme d'Ingénieur de Conception (Master 2) : \textbf{Mathématiques Appliquées}, Génie Logiciel, Algorithmique.}
\end{itemize}

% --- COMPÉTENCES ---
\section{Compétences Techniques}
\begin{itemize}[leftmargin=*, nosep]
    \item \small{\textbf{Backend \& API :} Java, Spring Boot / Spring WebFlux, PHP/Laravel, API REST, Architecture DDD, Microservices.}
    \item \small{\textbf{Messaging \& Data :} RabbitMQ, Redis (Cache/Pub-Sub), PostgreSQL, MySQL, MongoDB, ScyllaDB (NoSQL).}
    \item \small{\textbf{DevOps \& Outils :} Docker, Git, CI/CD, Tests unitaires, Agile/Scrum, Linux.}
    \item \small{\textbf{Langues :} Français (Natif), Anglais (Professionnel/Technique).}
\end{itemize}

% --- EXPÉRIENCES ---
\section{Expériences Professionnelles}
\begin{itemize}[leftmargin=*, label={-}]

\item \textbf{Stagiaire Recherche -- Sécurité Numérique \& NLP} \hfill \textit{Avril 2026 -- Août 2026} \\
\textit{INRIA Saclay -- Équipe TRIBE (Plateau de Saclay / IP Paris)} \hfill \textit{Python, NLP, BERT, pandas}
\vspace{-2pt}
    \begin{itemize}[leftmargin=0.4cm, nosep, label=\tiny$\bullet$]
        \item \small{Conception d'un pipeline Python d'analyse automatisée de la conformité RGPD d'applications mobiles.}
        \item \small{Analyse NLP à grande échelle des politiques de confidentialité avec modèles Transformers (BERT).}
        \item \small{Contribution à des publications scientifiques dans le cadre du programme national PEPR MOBIDEC.}
    \end{itemize}
\vspace{3pt}

\item \textbf{Développeur Backend -- Plateforme Bancaire Sécurisée} \hfill \textit{Oct. 2024 -- Jan. 2025} \\
\textit{CSC (Cameroon Software Company) -- Fintech} \hfill \textit{Laravel, MySQL, Redis, Docker}
\vspace{-2pt}
    \begin{itemize}[leftmargin=0.4cm, nosep, label=\tiny$\bullet$]
        \item \small{Conception du backend d'une plateforme de transactions financières sécurisée, conforme aux normes bancaires.}
        \item \small{Gestion de la concurrence des transactions et architecture scalable (Docker, Redis, tests unitaires).}
    \end{itemize}
\vspace{3pt}

\item \textbf{Développeur Full Stack -- Télémédecine} \hfill \textit{Fév. 2022 -- Mars 2024} \\
\textit{SOSDOCTA (Remote, Belgique/Canada)} \hfill \textit{Laravel, MySQL, OAuth2/JWT, API REST}
\vspace{-2pt}
    \begin{itemize}[leftmargin=0.4cm, nosep, label=\tiny$\bullet$]
        \item \small{Développement et maintenance (2 ans) d'une plateforme complète de télémédecine : API REST, réservation, paiement.}
        \item \small{Gestion de données médicales sensibles (conformité RGPD) et authentification sécurisée (OAuth2/JWT).}
    \end{itemize}
\vspace{3pt}

\end{itemize}

% --- PROJETS ---
\section{Projets Pertinents}
\begin{itemize}[leftmargin=*, label={-}]

\item \textbf{Snappy -- Plateforme de Messagerie Sécurisée (Projet académique)} \hfill \textit{2024 -- 2025} \\
\textit{Architecture DDD, 4 modules} \hfill \textit{Spring Boot / WebFlux, PostgreSQL, Redis, Socket.IO}
\vspace{-2pt}
    \begin{itemize}[leftmargin=0.4cm, nosep, label=\tiny$\bullet$]
        \item \small{Participation à la conception d'une architecture \textbf{Spring Boot réactive} (WebFlux, Mono/Flux) gérant des milliers de connexions WebSocket simultanées avec backpressure.}
        \item \small{Implémentation de cache multi-niveaux (Redis LRU + applicatif), indexation composite PostgreSQL, élimination des requêtes N+1.}
        \item \small{Chiffrement bout en bout (protocole Signal, AES-256-GCM) et conformité RGPD native.}
    \end{itemize}
\vspace{3pt}

\item \textbf{GoFind -- Architecture Microservices} \hfill \textit{2024} \\
\textit{Projet académique} \hfill \textit{Laravel, NestJS, MongoDB, RabbitMQ, Docker}
\vspace{-2pt}
    \begin{itemize}[leftmargin=0.4cm, nosep, label=\tiny$\bullet$]
        \item \small{Conception d'une architecture microservices complète avec communication asynchrone via \textbf{RabbitMQ} et API Gateway.}
        \item \small{Orchestration des services avec Docker et gestion de bases de données hétérogènes (MongoDB, MySQL).}
    \end{itemize}
\vspace{3pt}

\item \textbf{Laravel API Generator -- Outil Open Source} \hfill \textit{2024 -- Présent} \\
\textit{Projet personnel, +30 étoiles GitHub} \hfill \textit{PHP, Laravel, Python, LLM}
\vspace{-2pt}
    \begin{itemize}[leftmargin=0.4cm, nosep, label=\tiny$\bullet$]
        \item \small{Création d'un outil de génération automatique de code backend Laravel (modèles, contrôleurs, migrations, routes) à partir de diagrammes UML/XML. Déployé sur Packagist.}
    \end{itemize}
\vspace{3pt}

\end{itemize}

% --- DISTINCTIONS ---
\section{Distinctions \& Engagement}
\begin{itemize}[leftmargin=*, nosep]
    \item \small{\textbf{1\textsuperscript{ère} Place} -- IndabaX Africa 2025 (IA Santé : déploiement API, Flask, Streamlit). \textbf{1\textsuperscript{ère} Place} -- CITS National Hackathon (75 équipes).}
    \item \small{\textbf{Leadership :} Chef de la Cellule Projets et Cellule IA (Polytechnique Yaoundé). Animation workshops : Laravel, React, ML.}
    \item \small{\textbf{Certifications :} Supervised ML (DeepLearning.AI/Stanford), Python for Data Science (IBM), Jira (Coursera).}
\end{itemize}

\end{document}
